\documentclass[11pt,a4paper]{article}
\usepackage[T1]{fontenc}
\usepackage{lmodern,amsmath,amssymb,amsthm,microtype}
\usepackage[margin=26mm]{geometry}
\usepackage[colorlinks=true,linkcolor=blue,urlcolor=blue]{hyperref}
\newtheorem{lemma}{Lemma}[section]
\newtheorem{corollary}[lemma]{Corollary}
\newcommand{\Q}{\mathbb Q}
\title{Heights of specialized rational functions on a congruent curve}
\author{}\date{17 November 2026}
\begin{document}\maketitle

Fix $d\in\mathbf Z_{>0}$ and
$E_d:y^2=f(x)=x^3-d^2x$. All heights are absolute logarithmic
heights over $\Q$.

\section{A uniform specialization estimate}

Let $A,B,C,D\in\Q[Z]$ have degree at most $N\ge0$.
Pad their coefficient vectors with zeros and form the vectors
$\mathbf p,\mathbf q$ of $A(Z)+B(Z)y$ and $C(Z)+D(Z)y$
in the basis
\[
                  1,Z,\ldots,Z^N,y,Zy,\ldots,Z^Ny.
\]
Define
\[
 \Gamma=\sum_{i<j}(p_iq_j-p_jq_i)^2,\qquad
 H_c=h(1:\mathbf p:\mathbf q).
\]
For rational coefficients, $\Gamma\ne0$ is equivalent to
$\mathbf p,\mathbf q$ being linearly independent. In the function
field of $E_d$, it is therefore equivalent to the quotient
$(A+By)/(C+Dy)$ being defined and nonconstant.

\begin{lemma}\label{lem:coefficient}
Suppose $\Gamma\ne0$ and $U=(x_U,y_U)\in E_d(\Q)$ satisfies
$C(x_U)+D(x_U)y_U\ne0$. Put
\[
                   r=\frac{A(x_U)+B(x_U)y_U}{C(x_U)+D(x_U)y_U}.
\]
Then
\[
 h(r)\ge\frac12h(x_U)-H_c
               -\frac12\log\bigl(4(N+1)(d^2+2)\bigr).
\]
\end{lemma}
\begin{proof}
The polynomial
\[
  F_r(Z)=(A(Z)-rC(Z))^2-f(Z)(B(Z)-rD(Z))^2
\]
vanishes at $x_U$. It is not the zero polynomial. Indeed, if
$B-rD$ were nonzero and $F_r=0$, the squarefree cubic $f$
would be a square in $\Q(Z)$, which is impossible. If
$B-rD=0$, then $F_r=0$ also gives $A-rC=0$, contradicting
$\Gamma\ne0$.

Choose a common denominator $L$ for the coefficients of
$A,B,C,D$ such that
$(L,L\mathbf p,L\mathbf q)$ is a primitive integer vector.
Every entry has absolute value at most $\exp(H_c)$.
Write $r=a/b$ in lowest terms, $b>0$, so that
$\max(|a|,b)=\exp(h(r))$.
The coefficients of
$L(bA-aC)$ and $L(bB-aD)$ have absolute value at most
$2\exp(H_c+h(r))$. Hence the nonzero integer polynomial
\[
 L^2\bigl((bA-aC)^2-f(Z)(bB-aD)^2\bigr)
\]
has every coefficient bounded in absolute value by
\[
                 4(N+1)(d^2+2)\exp(2H_c+2h(r)).
\]

For a nonzero integer polynomial having a rational root
$a_0/b_0\ne0$ in lowest terms, the rational-root theorem
shows that $a_0$ divides its lowest nonzero coefficient and
$b_0$ divides its leading coefficient. Thus the height of the root is
at most the logarithm of the largest absolute coefficient.
The root zero has height zero and satisfies the same bound.
Apply this observation to $x_U$ and rearrange.
\end{proof}

\section{Coefficients depending on another point and an input}

Let $\mathcal R(X,H,U)$ be a fixed rational function on
$\mathbf A^1\times E_d\times E_d$, generically nonconstant
in $U$. Reduce it using $y_U^2=f(x_U)$, and write
\[
 \mathcal R=
 \frac{A(X,H,x_U)+B(X,H,x_U)y_U}
      {C(X,H,x_U)+D(X,H,x_U)y_U}.
\]
After clearing a fixed common coefficient denominator in
$X,H$, and then a fixed integer denominator, the coefficients
in $x_U$ are integer polynomials in $X,x_H,y_H$.
Retain the denominators used in this presentation as nonzero
guards. Let the partial degrees of these coefficient polynomials
in the three parameters
be bounded by $D_X,D_x,D_y$, respectively. Put
\[
                         K_H=2D_x+3D_y.
\]
The integer polynomials defining the coefficient vectors
give a uniform bound
\begin{equation}\label{eq:coeff-height}
 H_c\le D_Xh(X)+D_xh(x_H)+D_yh(y_H)+O_{\mathcal R,d}(1).
\end{equation}
For example, this follows by clearing the three rational
parameter denominators to their indicated partial degrees;
the remaining multiplicative constant is the maximum of one
and the coefficient $\ell^1$ norms of these finitely many
integer polynomials.

Let $\Gamma(X,H)$ be the sum of coefficient-minor squares
defined above. It is not identically zero on
$\mathbf A^1\times E_d$: this surface has Zariski-dense real
points, so a sum of rational-function squares can vanish
identically only when every summand vanishes.
Over rational specializations,
its zero locus is exactly the locus where all coefficient
minors vanish.

\begin{corollary}\label{cor:canonical}
There is a fixed constant $C_{\mathcal R,d}$ such that, at
every rational source point where the displayed denominators
and $\Gamma(X,H)$ are nonzero,
\[
 h(\mathcal R(X,H,U))
   \ge\widehat h(U)-K_H\widehat h(H)-D_Xh(X)-C_{\mathcal R,d}.
\]
\end{corollary}
\begin{proof}
The standard comparison gives
$|h(x_R)-2\widehat h(R)|\le C_d$ for nonzero rational
points $R$. Also the curve equation gives
\[
                 2h(y_H)\le3h(x_H)+\log(1+d^2).
\]
Insert these estimates into \eqref{eq:coeff-height} and
Lemma~\ref{lem:coefficient}. The canonical-height comparison
is as in \cite[Chapter VIII, Section9]{AEC}.
\end{proof}

In particular, suppose $U=[n]H$ for a nonzero integer $n$,
and let $\mathcal E$ be any quantity satisfying
\[
               \mathcal E\ge2h(X),\qquad
               \mathcal E\ge h(\mathcal R(X,H,U)).
\]
Then
\begin{equation}\label{eq:index-height}
 \mathcal E\ge\frac{2}{D_X+2}(n^2-K_H)\widehat h(H)-O_{\mathcal R,d}(1).
\end{equation}
This follows by multiplying the first inequality by $D_X/2$
and adding the second before applying the corollary.
The constants are uniform in the input and both points.

\begin{thebibliography}{9}
\bibitem{AEC}J.~H.~Silverman, \emph{The Arithmetic of Elliptic Curves},
second edition, Graduate Texts in Mathematics 106, Springer, 2009,
Chapter VIII, Section9.
\end{thebibliography}
\end{document}
